\documentclass[11pt]{article}
\usepackage[margin=1in]{geometry}
\usepackage{amsmath,amssymb,amsthm}
\usepackage[hidelinks]{hyperref}
\newtheorem{theorem}{Theorem}
\newtheorem{corollary}[theorem]{Corollary}
\begin{document}
\setlength{\emergencystretch}{2em}
\title{Ordinary binomial Strong Factorial through degree six}
\author{}\date{}\maketitle
Let $\mathcal L(X^j)=j!$. Here Strong Factorial for a nonzero
univariate phase $f$ means that
\[
 \bigl(\mathcal L(f^m),\mathcal L(f^{m+1})\bigr)\ne(0,0)
 \qquad\text{for every integer }m\ge1.
\]
This note assembles exact all-order family theorems and states the
next finite-support frontier. The assertions concern ordinary
univariate binomials; no reduction of arbitrary-support FC or SFC
to binomials is asserted.

\begin{theorem}[All ordinary binomials through degree six]
\label{thm:degree-six}
Every nonzero complex univariate binomial of degree at most six
satisfies Strong Factorial.
\end{theorem}
\begin{proof}
Corollary \texttt{cor:degree-six-frontier} of the
\href{QUINTIC_SCALAR_MINOR_SFC.tex}{quintic scalar-minor note}
reduces this assertion exactly to the supports
$\{2,6\},\{3,6\},\{4,6\}$ at terminal orders $N\ge6$.
Theorem \texttt{thm:fourth-gap} in
\href{LOWER_TWO_FOURTH_POWER_GAP_SFC.tex}{the boundary Frobenius note}
covers $X^2(A+BX^4)$ at every positive consecutive pair.
Theorem \texttt{thm:diagonal-prime} in
\href{DIAGONAL_PRIME_GAP_SFC.tex}{the diagonal-prime note}
at $p=3$ covers $X^3(A+BX^3)$ at every such pair.
Corollary \texttt{cor:ordinary-support-four-six} in
\href{LOWER_FOUR_QUADRATIC_GAP_SFC.tex}{the quadratic scalar-minor note}
covers $X^4(A+BX^2)$ at every such pair.
Thus every test in the exact reduction holds.
Each family includes the degenerate nonzero monomials.
\end{proof}

\begin{corollary}[The exact ordinary degree-seven frontier]
\label{cor:degree-seven-frontier}
Put
\[
 P_j^{a,d}(t)=\sum_{k=0}^j\binom jk(aj+1)_{dk}t^k.
\]
Strong Factorial for every ordinary univariate binomial of degree
at most seven is equivalent to
\[
 \gcd_{\mathbb Q[t]}(P_{N-1}^{a,d},P_N^{a,d})=1
\]
for every terminal integer $N\ge6$ and each of
\[
 (a,d)=(1,6),(3,4),(4,3),(5,2).
\]
These are precisely the supports
$\{1,7\},\{3,7\},\{4,7\},\{5,7\}$.
\end{corollary}
\begin{proof}
Theorem~\ref{thm:degree-six} covers every smaller degree.
For degree seven write the binomial as
$X^a(A+BX^{7-a})$, $0\le a\le6$, $AB\ne0$.
At $a=0$, Theorem \texttt{thm:all} of the
\href{ALL_UNIVARIATE_TWO_TERM_SFC.tex}{constant-plus-power note}
applies. At $a=2$, Theorem \texttt{thm:prime} of the
\href{LOWER_TWO_PRIME_GAP_SFC.tex}{lower-two prime-gap note}
applies with prime five and ordinary shape one.
At $a=6$, Theorem \texttt{thm:main} of the
\href{UNIFORM_PADE_PARABOLA_SFC.tex}{adjacent-exponent note}
applies. Exactly the four displayed pairs remain.

Theorem \texttt{thm:ordinary-fifth} of the
\href{lower_exponent_initial_and_boundary.tex}{initial-order note}
covers all ordinary binomials through terminal order five.
For every remaining pair,
\[
 \mathcal L\bigl((X^a(A+BX^d))^j\bigr)
       =A^j(aj)!P_j^{a,d}(B/A).
\]
This identifies the remaining tests with the displayed rational
gcds at $N\ge6$. Two rational polynomials have a common complex root
exactly when their rational gcd is nonconstant. A nonzero monomial
has every positive factorial moment nonzero.
Conversely each displayed support and order is among the degree-seven
tests. Thus the equivalence holds. It does not prove any of those
four unrestricted families.
\end{proof}

The input family theorems have written proofs, with their exact
algebraic reductions and internal replay records retained separately.
The present argument is a proof by specialization and support
classification. Independent mathematical review, formal verification
and priority are separate questions.
\end{document}
